% =============================================================
% intro_template.tex
% Five-paragraph introduction. Replace [BRACKETS] with content.
% Source via \input{tex/intro} from paper.tex.
% =============================================================

\section{Introduction}\label{sec:intro}

% --- 1. Motivation -------------------------------------------
[OPENING: substantive issue, not "This paper examines". One or two
sentences on why the question matters for economics, policy, or
welfare. End with a sharp question.]

% --- 2. What this paper does ---------------------------------
We study [QUESTION] in [SETTING]. Our analysis uses [DATA SOURCE,
SAMPLE PERIOD, UNIT OF OBSERVATION] and exploits [SOURCE OF
IDENTIFYING VARIATION].

% --- 3. Identification ---------------------------------------
The identifying assumption is that, absent [TREATMENT], [TREATED
UNITS] would have followed the same trend in [OUTCOME] as
[CONTROL UNITS]. We provide three pieces of supporting evidence:
(i) pre-period leads from an event-study specification are jointly
indistinguishable from zero; (ii) the same design applied to
[PLACEBO OUTCOME] yields a precise null; (iii) pre-treatment
characteristics are balanced across treatment cohorts. The
remaining identification concern is [HONEST CONCERN]; we address
it in \Cref{sec:robustness}.

% --- 4. Headline result --------------------------------------
We find that [TREATMENT] [increases / decreases] [OUTCOME] by
\headlineEstimate{} units (s.e. \headlineSE{}), or roughly
[X]\% relative to the sample mean. The effect is [LARGER /
CONCENTRATED IN / DRIVEN BY] [SUBGROUP], consistent with [MECHANISM].

% --- 5. Contribution to literature ---------------------------
Our paper contributes to three strands of literature. First, we
extend \citet{Author2020} by [SPECIFIC EXTENSION]. Second, we
provide new evidence on [MECHANISM] that complements
\citet{Other2019}. Third, we connect to the recent debate on
[BROADER QUESTION], where \citet{Third2022} argue [POSITION];
our findings [SUPPORT / QUALIFY / CHALLENGE] this view.

% --- (Optional) Road map ------------------------------------
% \Cref{sec:data} describes the data. \Cref{sec:strategy}
% details the empirical strategy. \Cref{sec:results} reports
% the main results, \Cref{sec:robustness} robustness checks,
% and \Cref{sec:conclusion} concludes.
